\section{Related work}\label{sec:related}

% P1. Supervised classification is competitive with or better than the other objectives at full
% downstream data, and no study varies its labels. Run-in heads are italic (no bold in the body).
\emph{Pretraining objectives for jets.} Supervised classification of simulated jets is the reference
objective for jet foundation models~\cite{part,sophon}, and the alternatives are measured against it.
The self-supervised families include masked particle modeling~\cite{mpm,mpmtok}, generative
pretraining~\cite{omnijet}, contrastive learning on resimulated jets~\cite{rs3l} and joint-embedding
predictive architectures~\cite{hepjepa,jpjepa}. Machine-learned particle flow, trained to reconstruct the
particles of an event, has been evaluated as a foundation model through its frozen features~\cite{mlpf_fm}. Other models combine classification with a second
objective: OmniLearn and OmniLearned with generation~\cite{omnilearn,omnilearned}, SuperJEPA with
joint-embedding prediction~\cite{superjepa}. One study compares the objectives on the same
pretraining jets, all from JetClass. It finds classification best or tied when the full top-tagging
training set is used for fine-tuning, with masked-particle pretraining ahead only for its smallest
model fine-tuned on \LitObjectivesFewJets{} jets, and concludes that the best objective is not a
property of the model alone~\cite[Sec.~IV.A]{objectives}. A comparison on JetClass-II evaluates the
objectives by fine-tuning alone, without linear probes~\cite{reprlearn}. Each of these studies keeps
the classes of its supervised objective fixed, or changes them together with the jets. The label set
of the objective that does best when labeled jets are plentiful has therefore not been varied on its
own.

% P2. Label granularity in vision: nearest precedent Chen Table VI; Hong's random per-image
% identities; Pirovano's grouping dependence; Mahajan linear versus fine-tuned; how our design differs.
\emph{Label granularity in vision.} The nearest precedent for our random partitions regroups whole
CIFAR-10 classes at random into two coarse classes. Training on the fine labels beats training on the
coarse ones for every grouping, by up to \LitChenRandomMaxGain{} percentage points of accuracy on the
coarse task~\cite[Table~VI]{chen}, and in the same study the gain from fine labels on CIFAR-100
becomes negligible once the coarse task has many classes~\cite[Table~VIII]{chen}. Both experiments
train and test on one task, so neither measures transfer. Whether fine labels help also depends on
how the classes are grouped. For cats and dogs against cars and trucks, training on the four fine
classes never helps, while for a grouping that mixes animals and vehicles it consistently
does~\cite[Fig.~3]{pirovano}. Class identities assigned at random to individual images, rather than
to whole classes, do worse than every level of a species taxonomy~\cite{hong}. With a linear probe on
frozen features, pretraining on a larger set of web hashtags leads by about
\LitMahajanLinearGapPercent{} percentage points of accuracy on the ImageNet task with the most
classes, and under full fine-tuning the lead nearly disappears~\cite{mahajan}.\footnote{After full
fine-tuning, the larger label set leads on that task by a margin comparable to the difference the
same study finds between two pretrainings on different random samples of images. The smaller label
set, whose hashtags match the ImageNet classes, leads on the standard ImageNet task, and it was also
pretrained on fewer images~\cite[Fig.~1, Sec.~3.1]{mahajan}.} With the labels fixed, losses that
separate the training classes more give frozen features from which a linear probe transfers worse,
while after fine-tuning the losses are indistinguishable~\cite{kornblith}. Our random partitions
differ from these controls in three ways. They keep the number of classes of a physical level,
\Kcoarse{} groups; they match each group's share of the training jets to that level; and they are
pretraining label sets evaluated on fixed downstream tasks, not targets in their own right.

% P3. Auxiliary physics targets: GN2, GloParT (argument, not measurement), CMS, GN3X; none measures
% the effect on frozen features with the labels fixed.
\emph{Auxiliary physics targets.} Production taggers add auxiliary targets to the classification,
with effects that point both ways. The ATLAS flavor tagger GN2 has two auxiliary objectives, which
predict the origin of each track and which tracks share a vertex. Removing both lowers its $c$- and
light-jet rejection by up to \LitGnTwoAuxRejectionDropPercent\%, and keeping either one recovers most
of the loss~\cite[Discussion]{gntwo}. GloParT adds two log-cosh regressions of the jet mass to its
classification loss, with a weight tuned to the value our mass models also use. Its author argues
that because the terms do not interact, the minimum of the combined loss is optimal for both
tasks~\cite[Sec.~11.2.1]{glopart}; no measurement accompanies the argument. A CMS tagger of boosted
Higgs bosons decaying to pairs of $W$ bosons regresses the mass in GeV, so its weight is on a
different scale from ours~\cite{cmspart}. ATLAS measured a cost instead: a network for boosted Higgs
bosons trained for classification and mass regression together was about
\LitAtlasUnifiedEffDropPercent\% less efficient at classification than a dedicated classifier, and
ATLAS regresses the mass in a separate network~\cite[Sec.~6.5]{atlasgnx}. These studies measure the
tagger's outputs. None measures what an auxiliary target does to the frozen features while the labels
stay fixed, which our \Kres- and \Kcoarse-class models, pretrained with and without the mass
regression, do.

% P4. Transfer from frozen features: Sophon's output discriminant and transfer head; Vigl's
% fine-tuned versus frozen; neither varies the label set; Sophon reports no variation between runs.
\emph{Transfer from frozen features.} Sophon used its label set downstream in two
ways~\cite[Sec.~III]{sophon}. Its output scores form an $X\to b\bar b$ discriminant directly. A small
network trained on its frozen features, with its first layer taken from the classification head,
identifies $X\to bs$ decays, which no Sophon class contains. At low background efficiency this
network does better than the Particle Transformer and ParticleNet trained on the task~\cite[Fig.~2,
App.~B.5]{sophon}. Vigl et~al.\ compared frozen and fine-tuned use of a Particle Transformer backbone
pretrained on $X\to b\bar b$ jets in a search for $G\to HH\to b\bar b b\bar b$, and fine-tuning was
better for every architecture they tried~\cite{vigl}. Neither study varied the label set with the
training jets fixed, and Sophon reports no variation between pretraining runs. We repeat Sophon's % CHECK: pending Xbx* (drop this sentence with section 4.7 if the samples are not staged by 12 Nov)
benchmark at every level of the hierarchy.
